\begin{annexitem}{annex:p03.cv.1}{thm:p03-cv-support}{Compactness bounds semantic support}
\begin{proof}
Fix \(\varphi\in\mathscr L[\tau]\). Choose a disjoint renamed
copy \(\rho\tau\), including disjoint copies of the sorts.
For each original sort \(s\), add a binary relation \(B_s\)
of argument sorts \(s,\rho(s)\), and first-order axioms saying
that it is a bijection between those domains. Add compatibility
axioms for each non-sort symbol. For a relation \(R\), they say
\[
 \bigwedge_i B_{s_i}(x_i,y_i)
       \ \longrightarrow\
       \bigl(R\bar x\leftrightarrow\rho(R)\bar y\bigr),
\]
with universal closure and the appropriate argument sorts. For a
function \(f\) of output sort \(t\), use
\[
 \begin{aligned}
 &\left(\bigwedge_i B_{s_i}(x_i,y_i)\land B_t(z,w)\right)\\
 &\hspace{2em}\longrightarrow
       \bigl(f(\bar x)=z\leftrightarrow\rho(f)(\bar y)=w\bigr).
 \end{aligned}
\]
For a constant \(c\) of sort \(t\), use
\[
 B_t(x,y)\ \longrightarrow\ (x=c\leftrightarrow y=\rho(c)).
\]
Nullary relations use the same clause with an empty antecedent.
Let \(\Delta\) be this bridge theory. All its sentences are
first-order and hence available from the stated closure properties.
There are at most \(\kappa\) of them, because
\(|\tau|\leq\kappa\), all arities are finite, and \(\kappa\)
is infinite.

In any model of \(\Delta\), the bridges give a sorted
isomorphism from the original reduct to the renamed reduct with
its names changed back. Isomorphism invariance, renaming and
reduct invariance therefore give
\(\varphi\leftrightarrow\varphi^\rho\).
Thus \(\Delta\cup\{\neg(\varphi\leftrightarrow\varphi^\rho)\}\)
has no model. Its size is at most \(\kappa\).
By \((\kappa,\lambda)\)-compactness, an unsatisfiable subset
has size less than \(\lambda\). Remove the displayed negated
biconditional from that subset, and call the remaining bridge
axioms \(\Delta_0\). Then
\[
 \Delta_0\models\varphi\leftrightarrow\varphi^\rho,
 \qquad |\Delta_0|<\lambda.
\]
Every sentence of \(\Delta_0\) mentions finitely many original
symbols and sorts, their copies, and the corresponding bridges.
Let \(\sigma\) contain those original symbols and all their
attached sorts. If the collection is empty, include one original
sort to obtain a vocabulary. Its size is less than \(\lambda\):
a finite collection of finite vocabularies is finite, and an
infinite collection of size \(\mu<\lambda\) has union of size
at most \(\mu\). No regularity assumption on the cardinal
\(\lambda\) is used.

Suppose \(\mathfrak A,\mathfrak B\) have isomorphic
\(\sigma\)-reducts. Build a structure interpreting the original
sorts and symbols by \(\mathfrak A\), and their disjoint copies
by \(\mathfrak B\). On every sort mentioned by \(\Delta_0\),
interpret \(B_s\) by the given reduct isomorphism. Interpret
unused bridges arbitrarily, for example by the empty relation.
The selected bijection axioms hold, and each selected symbol
compatibility holds because its symbol and attached sorts are
in \(\sigma\). Hence this structure satisfies \(\Delta_0\).
The displayed consequence and the renaming property give
\(\mathfrak A\models\varphi\) if and only if
\(\mathfrak B\models\varphi\). Thus \(\sigma\) is a semantic
support. This argument does not assume that unselected sorts have
equal cardinalities or that the two original universes are equal.

For a compact logic choose an infinite \(\kappa\geq|\tau|\)
and put \(\lambda=\aleph_0\). Full compactness implies the
required cardinal compactness, and the resulting support is finite.
\end{proof}
\end{annexitem}

\begin{annexitem}{annex:p03.cv.2}{thm:p03-cv-support-translation}{A sentence over its semantic support}
\begin{proof}
Retain the single sort, choose a fresh constant \(d\), and
canonically interpret every symbol removed from \(\tau\): each
relation is false, each function has constant value \(d\), and
each removed constant is \(d\). The function graph templates are
\(y=d\), and the constant templates are the same equality.
They are total and single-valued in every nonempty structure with
an interpretation of \(d\).

Admissible simultaneous substitution gives
\[
 \varphi^*(d)\in\mathscr L[\sigma\cup\{d\}].
\]
For a \(\sigma\)-structure and any value of \(d\), its truth is
exactly the truth of \(\varphi\) in this canonical expansion.
Every other \(\tau\)-expansion has the same \(\sigma\)-reduct,
so semantic support makes that truth independent of all the
removed interpretations and of the chosen value of \(d\).
Particularization gives \(\psi=\exists d\,\varphi^*(d)\) in
\(\mathscr L[\sigma]\). Nonemptiness provides a possible value
of \(d\), and the preceding equivalences show that \(\psi\)
is equivalent to \(\varphi\) in every \(\tau\)-structure.
Combining this with \cref{thm:p03-cv-support} proves the compact
case. No elimination of sorts is inferred from the one-sorted
substitution operation.
\end{proof}
\end{annexitem}

\begin{annexitem}{annex:p03.cv.3}{thm:p03-cv-space}{Logical and topological compactness}
\begin{proof}
Actual complete theories satisfy
\[
 \begin{aligned}
 U_\varphi\cap U_\psi&=U_{\varphi\land\psi},\\
 X_{\mathscr L,\tau}\setminus U_\varphi&=U_{\neg\varphi}.
 \end{aligned}
\]
A tautology supplies the whole space. Hence these sets form a
clopen basis. Distinct points differ on a sentence; that sentence
and its negation give disjoint basic neighbourhoods. The space
is therefore Hausdorff.

Assume the local logical compactness property. Given an open cover,
refine it by taking all basic opens contained in one of its members.
If this basic cover had no finite subcover, the negations of its
defining sentences would be finitely satisfiable. A model of all
those negations would give a point of the theory space outside
the cover, a contradiction. Thus finitely many basic opens cover;
selecting a containing original open for each gives a finite
subcover of the original cover.

Conversely, let \(T\subseteq\mathscr L[\tau]\) be finitely
satisfiable. The closed sets \(U_\varphi\), for \(\varphi\in T\),
have the finite intersection property. Compactness of the space
gives a point in their total intersection. By definition the point
is the theory of an actual structure, which consequently satisfies
\(T\). Applying this argument separately to every vocabulary
proves the global assertion for small logics.

For the representative-structure presentation of the source, the
map sending a chosen representative to its complete theory is a
bijection and carries each sentence-defined basic open to
\(U_\varphi\). It and its inverse are continuous by the basis
criterion, so the two presentations are homeomorphic.
\end{proof}
\end{annexitem}

\begin{annexitem}{annex:p03.cv.4}{thm:p03-cv-blocks}{Closed sets and theory blocks}
\begin{proof}
For a theory \(T_i\), put
\(F_i=\bigcap_{\varphi\in T_i}U_\varphi\), a closed set.
For any \(J\subseteq I\), its intersection over \(i\in J\)
is nonempty exactly when \(\bigcup_{i\in J}T_i\) has a model:
a point is precisely the complete theory of such a model.
Thus topological \((\kappa,\lambda)\)-compactness implies the
block condition.

Every closed set is of this form. Indeed its open complement is
a union of basic opens; negating the sentences defining those
opens expresses the closed set as their closed intersection.
Consequently the block condition, applied to theories defining
a family of closed sets, proves the converse as well.

Ordinary logical \((\kappa,\lambda)\)-compactness is the
special case where the blocks are singleton theories, so it
concerns basic clopens defined by individual sentences. It does
not in general bound the total number of sentences in arbitrary
blocks. If \(|\mathscr L[\tau]|\leq\kappa\), however, the
whole union of any blocks has size at most \(\kappa\).
Any subset of that union of size less than \(\lambda\) uses
sentences from fewer than \(\lambda\) blocks: choose one index
for each sentence. The hypothesis of the block condition makes
that subset satisfiable. Logical \((\kappa,\lambda)\)-compactness
now gives a model of the whole union, proving the last assertion.
\end{proof}
\end{annexitem}

\begin{annexitem}{annex:p03.cv.5}{thm:p03-cv-qone-space}{A compactness boundary}
\begin{proof}
Countable compactness of \(\operatorname{FO}(Q_1)\) is
\cref{thm:p02-w02c-p02-003-keisler-compactness}. If formulated
in countable vocabulary, it applies to an arbitrary countable
theory by taking the countable union of the finite vocabularies
of its sentences, then expanding unused symbols arbitrarily.

By Cantor's theorem and Choice, there is an injection
\(\alpha\mapsto e_\alpha\) from \(\omega_1\) into
\(\{0,1\}^{\omega}\). Let \(\tau\) contain a constant
\(c_\alpha\) for each \(\alpha<\omega_1\), and let
\(\chi=\neg Q_1x\,(x=x)\), the assertion that the domain
is at most countable. For \(n<\omega\), put
\[
 T_n=\{\chi\}\cup
     \{c_\alpha\ne c_\beta:
                     e_\alpha(n)\ne e_\beta(n)\}.
\]
The language has a set of sentences, so its theory space is defined.
Each \(T_n\) determines the closed set
\(F_n=\bigcap_{\varphi\in T_n}U_\varphi\).

For a finite \(J\subseteq\omega\), interpret \(c_\alpha\)
by the binary vector \(e_\alpha|J\), and take as domain the
nonempty finite set of vectors so obtained. Distinct coordinate
values at any \(n\in J\) force distinct vectors, so every
inequality from the corresponding \(T_n\) holds. The domain
is finite and hence satisfies \(\chi\). Thus
\(\bigcup_{n\in J}T_n\) has a finite model, and
\(\bigcap_{n\in J}F_n\ne\varnothing\).

The full union has no model. For \(\alpha\ne\beta\), the
binary sequences differ at some coordinate \(n\), so that union
contains \(c_\alpha\ne c_\beta\). A model would therefore
contain \(\aleph_1\) distinct elements, while also satisfying
\(\chi\). Hence \(\bigcap_{n<\omega}F_n=\varnothing\).
This is a countable family of closed sets with the finite
intersection property but empty intersection, proving the
failure of topological countable compactness. The blocks can
be uncountable; they are not single sentences to which logical
countable compactness could be applied. Only the existence of
uncountably many distinct binary sequences was used, not the
continuum hypothesis.
\end{proof}
\end{annexitem}
